%Data institusi

%\input{Dozentdaten}

%\renewcommand{\fachbereich}{Fachbereich}

%\renewcommand{\dozent}{Prof. Dr. . }

%Data ujian

\renewcommand{\klausurgebiet}{ }

\renewcommand{\klausurtyp}{ }

\renewcommand{\klausurdatum}{ . 20}

\klausurvorspann {\fachbereich} {\klausurdatum} {\dozent} {\klausurgebiet} {\klausurtyp}

%Data untuk tabel nilai berikut

\renewcommand{\aeins}{  3  }

\renewcommand{\azwei}{  3   }

\renewcommand{\adrei}{  0  }

\renewcommand{\avier}{  4  }

\renewcommand{\afuenf}{  5  }

\renewcommand{\asechs}{  4  }

\renewcommand{\asieben}{  9  }

\renewcommand{\aacht}{  0  }

\renewcommand{\aneun}{  8  }

\renewcommand{\azehn}{  6  }

\renewcommand{\aelf}{  0  }

\renewcommand{\azwoelf}{  3  }

\renewcommand{\adreizehn}{  3   }

\renewcommand{\avierzehn}{  0  }

\renewcommand{\afuenfzehn}{  0  }

\renewcommand{\asechzehn}{  7  }

\renewcommand{\asiebzehn}{  55  }

\renewcommand{\aachtzehn}{    }

\renewcommand{\aneunzehn}{    }

\renewcommand{\azwanzig}{    }

\renewcommand{\aeinundzwanzig}{    }

\renewcommand{\azweiundzwanzig}{    }

\renewcommand{\adreiundzwanzig}{    }

\renewcommand{\avierundzwanzig}{    }

\renewcommand{\afuenfundzwanzig}{    }

\renewcommand{\asechsundzwanzig}{    }

\punktetabellesechzehn

\klausurnote

\newpage

\setcounter{section}{K}

\inputaufgabegibtloesung
{3}
{

Definisikan istilah-istilah berikut
\zusatzklammer {yang dicetak miring} {} {.}
\aufzaehlungsechs{
\stichwort {Kelengkungan Gauss} {}
dari suatu permukaan diferensiabel

\mavergleichskette
{\vergleichskette
{ Y
}
{ \subseteq }{ \R^3
}
{  }{
}
{    }{
}
{   }{
}
}
{}{}{}
di suatu titik

\mavergleichskette
{\vergleichskette
{ P
}
{ \in }{ Y
}
{  }{
}
{    }{
}
{   }{
}
}
{}{}{.}

}{Suatu
\stichwort {manifold topologis} {}
berdimensi $n$.

}{Suatu
\stichwort {medan vektor tak bergantung waktu} {}
pada manifold diferensiabel $M$.

}{Suatu
\stichwort {bentuk diferensial berderajat $k$} {}
pada manifold diferensiabel $M$.

}{Suatu \stichwort {manifold Riemann} {.}

}{Suatu koneksi linear yang
\stichwort {bebas torsi} {}
pada bundel tangen manifold Riemann $M$.
}

}
{} {}

\inputaufgabegibtloesung
{3}
{

Nyatakan teorema-teorema berikut.
\aufzaehlungdrei{
\stichwort {Teorema eksistensi transpor paralel pada hiperpermukaan} {.}}{Teorema tentang deskripsi lokal bentuk diferensial.}{
\stichwort {Versi balok} {}
dari teorema Stokes.}

}
{} {}

\inputaufgabe
{0}
{

}
{} {}

\inputaufgabegibtloesung
{4 (1+3)}
{

Misalkan $M$ elips yang diberikan oleh

\mavergleichskettedisp
{\vergleichskette
{ x^2+3y^2
}
{ =} { 2
}
{ } {
}
{ } {
}
{ } {
}
}
{}{}{}
dalam $\R^2$,

\mavergleichskette
{\vergleichskette
{ P
}
{ = }{ \begin{pmatrix}  1  \\ { \frac{   1  }{  \sqrt{3}  } }    \end{pmatrix}
}
{  }{
}
{    }{
}
{   }{
}
}
{}{}{}
dan

\mavergleichskette
{\vergleichskette
{ Q
}
{ = }{ \begin{pmatrix}  -1 \\ { \frac{   1  }{  \sqrt{3}  } }    \end{pmatrix}
}
{  }{
}
{    }{
}
{   }{
}
}
{}{}{}
adalah titik-titik pada $M$.
\aufzaehlungdrei{Tunjukkan bahwa

\mavergleichskettedisp
{\vergleichskette
{ v
}
{ =} { \begin{pmatrix}  -  \sqrt{3}  \\ 1    \end{pmatrix}
}
{ } {
}
{ } {
}
{ } {
}
}
{}{}{}
merupakan
\definitionsverweis {vektor tangen}{}{}
di $P$ pada $M$.
}{Tentukan citra $v$ di bawah suatu transpor paralel pada $M$ dari $P$ ke $Q$.
}{
}

}
{} {}

\inputaufgabegibtloesung
{5}
{

Buktikan teorema tentang orientasi-orientasi pada hiperpermukaan diferensiabel.

}
{} {}

\inputaufgabegibtloesung
{4}
{

Misalkan

\mavergleichskette
{\vergleichskette
{ Y
}
{ \subseteq }{ W
}
{ \subseteq }{ \R^3
}
{    }{
}
{   }{
}
}
{}{}{,}
$W$
\definitionsverweis {terbuka}{}{,}
dan suatu
\definitionsverweis {permukaan diferensiabel}{}{}
yang dilengkapi dengan
\definitionsverweis {medan normal satuan}{}{}
$N$. Misalkan
\maabb {\gamma} { I } { Y
} {}
suatu kurva terdiferensialkan secara kontinu dan
\maabbdisp {F} { I } { \R^3
} {}
suatu
\definitionsverweis {medan vektor tangensial}{}{}
diferensiabel sepanjang $\gamma$ yang
\definitionsverweis {paralel}{}{.}
Tunjukkan bahwa medan vektor
\maabbeledisp {} { I } { \R^3
} { t } { N(\gamma(t)) \times F(t)
} {,}
juga paralel, dengan $\times$ menyatakan
\definitionsverweis {hasil kali silang}{}{}
dalam $\R^3$.

}
{} {}

\inputaufgabegibtloesung
{9}
{

Misalkan

\mavergleichskette
{\vergleichskette
{ M
}
{ \neq }{ \emptyset
}
{  }{
}
{    }{
}
{   }{
}
}
{}{}{}
suatu
\definitionsverweis {manifold diferensiabel}{}{}
berdimensi $n$. Tunjukkan bahwa terdapat suatu rantai
\definitionsverweis {submanifold-submanifold tertutup}{}{}

\mavergleichskettedisp
{\vergleichskette
{ M_0
}
{ \subseteq} { M_1
}
{ \subseteq} { M_2
}
{ \subseteq  \ldots \subseteq} { M_{n-1}
}
{ \subseteq} { M_n
}
}
{
\vergleichskettefortsetzung
{ =} { M
}
{ } {}
{ } {}
{ } {}
}{}{}
sedemikian sehingga submanifold tertutup $M_i$ berdimensi $i$.

}
{} {}

\inputaufgabe
{0}
{

}
{} {}

\inputaufgabegibtloesung
{8}
{

Misalkan
\maabbdisp {\psi} {\R^n } {\R
} {}
suatu
\definitionsverweis {fungsi diferensiabel}{}{}
dan

\mavergleichskettedisphandlinks
{\vergleichskettedisphandlinks
{M
}
{ =} { \left\{  \left(  x_1   , \,   , \ldots ,  , \,   x_n , \,   \psi(x_1 , \ldots , x_n)                \right)   \mid  (x_1 , \ldots , x_n) \in \R^n         \right\}
}
{ \subseteq} { \R^n \times \R
}
{ } {
}
{ } {
}
}
{}{}{}
merupakan
\definitionsverweis {grafik}{}{}
dari $\psi$. Tunjukkan bahwa $M$ suatu
\definitionsverweis {manifold Riemann}{}{}
yang
\definitionsverweis {berorientasi}{}{}
dan bahwa untuk
\definitionsverweis {bentuk volume kanonik}{}{}
$\omega$ pada $M$ berlaku rumus

\mavergleichskettedisp
{\vergleichskette
{ (\operatorname{Id} \times \psi)^*(\omega)
}
{ =} { \left(  1+ \sum_{i = 1}^n \left(  { \frac{   \partial   \psi  }{  \partial   x_i   } }  \right)^2  \right)^{1/2} dx_1 \wedge \ldots \wedge dx_n
}
{ } {
}
{ } {
}
{ } {
}
}
{}{}{.}

}
{} {}

\inputaufgabe
{6}
{

Hitung luas permukaan bola satuan.

}
{} {}

\inputaufgabe
{0}
{

}
{} {}

\inputaufgabegibtloesung
{3}
{

Hitung
\definitionsverweis {turunan eksterior}{}{}
$d \omega$ dari
$1$-\definitionsverweis {bentuk diferensial}{}{}

\mavergleichskettedisp
{\vergleichskette
{ \omega
}
{ =} { { \frac{   x^2 }{  y } } dx- { \frac{   x  }{  y^2 } } dy
}
{ } {
}
{ } {
}
{ } {
}
}
{}{}{}
pada

\mavergleichskette
{\vergleichskette
{ U
}
{ = }{ \left\{ (x,y) \in \R^2  \mid   y \neq 0         \right\}
}
{  }{
}
{    }{
}
{   }{
}
}
{}{}{.}

}
{} {}

\inputaufgabegibtloesung
{3}
{

Berikan contoh suatu
\definitionsverweis {manifold berbatas}{}{}
yang
\definitionsverweis {terhubung}{}{,}
dan batasnya terdiri atas
\zusatzklammer {gabungan saling lepas dari} {} {}
tak hingga banyak lingkaran $S^1$.

}
{} {}

\inputaufgabe
{0}
{

}
{} {}

\inputaufgabe
{0}
{

}
{} {}

\inputaufgabegibtloesung
{7}
{

Misalkan $M$ suatu
\definitionsverweis {manifold Riemann}{}{}
dan diberikan suatu
\definitionsverweis {koneksi linear}{}{}
pada
\definitionsverweis {bundel tangen}{}{}
$TM$ yang
\definitionsverweis {metrik}{}{}
dan
\definitionsverweis {bebas torsi}{}{.}
Tunjukkan bahwa untuk medan-medan vektor $X,Y,Z$ berlaku rumus

\mavergleichskettedisp
{\vergleichskette
{ \left\langle \nabla_X Y , Z  \right\rangle
}
{ =} { { \frac{   1  }{  2 } }  \left(  D_X \left\langle  Y  ,  Z  \right\rangle +  D_Y \left\langle  Z  ,  X  \right\rangle -  D_Z \left\langle  X  ,  Y  \right\rangle  - \left\langle  Y  , [X,Z]   \right\rangle - \left\langle  Z  , [Y,X]   \right\rangle + \left\langle  X  , [Z,Y]   \right\rangle   \right)
}
{ } {
}
{ } {
}
{ } {
}
}
{}{}{.}

}
{} {}
